\documentclass[12pt,a4paper]{article}

% Idioma, codificación y tipografía
\usepackage[utf8]{inputenc}
\usepackage[T1]{fontenc}
\usepackage[spanish,es-nodecimaldot]{babel}
\usepackage{lmodern}

% Página, gráficos y tablas
\usepackage[margin=2.5cm]{geometry}
\usepackage{graphicx}
\usepackage[table]{xcolor} % Permite pintar celdas de gris
\usepackage{array}
\usepackage{booktabs}
\usepackage{multirow}
\usepackage{tabularx}
\usepackage{longtable}
\usepackage{enumitem}
\usepackage{fancyhdr}
\usepackage{hyperref}

\hypersetup{
    colorlinks=true,
    linkcolor=black,
    urlcolor=blue,
    pdftitle={Usability Test Dashboard 2.0 -- Sprint 1},
    pdfauthor={Equipo del proyecto}
}

\setlength{\parindent}{0pt}
\setlength{\parskip}{0.55em}
\renewcommand{\arraystretch}{1.25}
\setlist[itemize]{leftmargin=*,nosep}
\setlist[enumerate]{leftmargin=*,nosep}

\pagestyle{fancy}
\fancyhf{}
\fancyhead[C]{%
    \begin{minipage}[c]{1.5cm}
        \centering
        \IfFileExists{imagenes/uta.png}{%
            \includegraphics[width=1.5cm]{imagenes/uta.png}%
        }{%
            \makebox[1.5cm][c]{\rule{0pt}{1.5cm}}%
        }
    \end{minipage}%
    \hfill
    \begin{minipage}[c]{10cm}
        \centering
        \small {\textbf{UNIVERSIDAD TÉCNICA DE AMBATO}} \\[0.1cm]
        \scriptsize \textit{FACULTAD DE INGENIERÍA EN SISTEMAS ELECTRÓNICA E INDUSTRIAL} \\[0.1cm]
        \scriptsize \textbf{CARRERA SOFTWARE} \\[0.1cm]
        \scriptsize \textbf{CICLO ACADÉMICO: JULIO - DICIEMBRE 2026}
        \vspace{0.2cm}
    \end{minipage}%
    \hfill
    \begin{minipage}[c]{1.5cm}
        \centering
        \IfFileExists{imagenes/fisei.png}{%
            \includegraphics[width=1.5cm]{imagenes/fisei.png}%
        }{%
            \makebox[1.5cm][c]{\rule{0pt}{1.5cm}}%
        }
    \end{minipage}%
}
\fancyfoot[C]{\thepage}
\setlength{\headheight}{58pt}
\setlength{\headsep}{12pt}

\newcolumntype{Y}{>{\raggedright\arraybackslash}X}
\newcolumntype{L}[1]{>{\raggedright\arraybackslash}p{#1}}

\begin{document}

\thispagestyle{fancy}
\begin{center}
    {\Large\bfseries DOCUMENTACIÓN DEL SPRINT 1\par}
\end{center}

\section{Portada}

\begin{tabularx}{\textwidth}{@{}>{\bfseries\raggedright\arraybackslash}p{6.2cm}X@{}}
    Tema: & Usability Test Dashboard 2.0 -- Sprint 1 \\
    Unidad de Organización Curricular: & PROFESIONAL \\
    Nivel y Paralelo: & 5 - A \\
    Alumnos participantes: &  Cholota Bastidas Ariel Isai \\
                                             & García Amores Manolo José \\
                                             & Miranda Tarupi Luis Sebastián \\
                                             & Velasco Huayamave Kevin Andrés \\
    Asignatura: & Interacción Humano Computador (IHC) \\
    Docente: & Ing. Caiza Caizabuano José Ruben\\
\end{tabularx}

\vspace{1cm}
\tableofcontents
\clearpage
\pagenumbering{arabic}

\section{Contexto y objetivo del proyecto}

\textit{Usability Test Dashboard 2.0} es una plataforma para que evaluadores UX
incorporen imágenes o bocetos de interfaces y observaciones de pruebas de usabilidad.
El sistema utilizará inteligencia artificial para analizar esa información y generar
historias de usuario para el Product Backlog.

\subsection{Objetivo del Sprint 1}

Establecer la base del producto y avanzar en la carga de bocetos y en la definición
de la experiencia y la arquitectura necesarias para cumplir el objetivo de respuesta
de la inteligencia artificial. El sprint comprende del 23 al 30 de septiembre de 2026,
con una capacidad total planificada de 32 horas para el equipo.

\subsection{Definition of Done del Sprint 1}

Para considerar completo el Sprint 1, se recopilarán las evidencias que aplican a
esta iteración. La matriz solicita que los artefactos visuales se entreguen en formato
PNG.

\begin{table}[htbp]
    \centering
    \caption{Definition of Done y evidencias requeridas para el Sprint 1}
    \label{tab:definition-of-done}
    \begin{tabularx}{\textwidth}{L{2.7cm}Y L{2.7cm} L{2.0cm}}
        \toprule
        \textbf{Área} & \textbf{Evidencia requerida} & \textbf{Herramienta} & \textbf{Formato / estado} \\
        \midrule
        UX Research & User Persona y Journey Map & Miro o FigJam & PNG; pendiente \\
        Diseño UI/UX & Mockups que muestren jerarquía visual & Figma & PNG; pendiente \\
        Arquitectura de información & User Flow & Figma o Miro & PNG; pendiente \\
        Agile UX & Sprint Backlog gestionado en tablero Kanban & Trello & PNG; tablero estructurado, captura pendiente \\
        GitHub & Repositorio funcional, código fuente y documentación & GitHub & PNG de evidencia; repositorio creado, validar contenido \\
        Evidencia técnica & Proyecto base organizado y funcional & VS Code & PNG; pendiente \\
        \bottomrule
    \end{tabularx}
\end{table}

La Definition of Done anterior corresponde únicamente al Sprint 1. Los criterios de
la matriz destinados a sprints futuros quedan fuera del alcance de este documento.

\section{Equipo y responsabilidades}

\begin{table}[htbp]
    \centering
    \caption{Roles oficiales y responsabilidades del equipo}
    \label{tab:roles}
    \begin{tabularx}{\textwidth}{L{2.6cm} L{2.0cm}Y}
        \toprule
        \textbf{Integrante} & \textbf{Rol} & \textbf{Responsabilidades principales} \\
        \midrule
        Luis & Product Owner & Priorizar el Product Backlog, definir historias de usuario y validar que el asistente evalúe la usabilidad. Apoya en pruebas de Figma, validación UX y diseño de prompts para IA. \\
        Kevin & Scrum Master & Facilitar Daily Standups y planificación, gestionar Trello y eliminar bloqueos. Apoya en GitHub Projects, DevOps básico, Docker y desarrollo. \\
        Manolo & Development Team & Participar en arquitectura, lógica de negocio, API, base de datos e integración de IA. \\
        Ariel & Development Team & Participar en arquitectura, lógica de negocio, frontend, base de datos e integración de IA. \\
        \bottomrule
    \end{tabularx}
\end{table}

\section{Análisis de capacidad}

La capacidad disponible comunicada para el Sprint 1 es de 32 horas de equipo. La
distribución de esas horas entre integrantes y actividades debe mantenerse en el
tablero de Trello para hacer visible el trabajo comprometido y el trabajo restante.
Los puntos de historia incluidos en este documento son estimaciones sugeridas y no
han sido confirmados por el Development Team; no representan una conversión directa
de las 32 horas disponibles.

\section{Sprint Backlog}

Las historias de usuario y el requerimiento no funcional en progreso para este sprint
son los siguientes. Los puntos y responsables que aparecen en la tabla son propuestas
iniciales: deben validarse con el Development Team y registrarse en Trello antes de
confirmar el compromiso del Sprint Backlog.

\begin{table}[htbp]
    \centering
    \caption{Elementos en progreso del Sprint 1}
    \label{tab:sprint-backlog}
    \begin{tabularx}{\textwidth}{L{1.5cm}Y L{2.3cm} L{1.7cm} L{1.8cm}}
        \toprule
        \textbf{ID} & \textbf{Elemento} & \textbf{Responsable propuesto} & \textbf{Puntos propuestos} & \textbf{Estado} \\
        \midrule
        RF-01 & Subida de bocetos & Manolo y Ariel; Luis valida & 5 SP (provisional) & En progreso \\
        RNF-01 & Tiempo de respuesta IA menor a 10 s & Manolo y Ariel; Luis valida UX & 8 SP (provisional) & En progreso \\
        \bottomrule
    \end{tabularx}
\end{table}

\subsection{Trabajo previsto por elemento}

\begin{itemize}
    \item \textbf{RF-01:} análisis de requerimientos base, definición del flujo de usuario y diseño de las pantallas de carga.
    \item \textbf{RNF-01:} investigación de soluciones UX para gestionar la latencia (indicadores de carga y esqueletos) y definición del stack backend.
\end{itemize}

\section{Análisis y requerimientos}

\subsection{RF-01: Requerimientos funcionales base, módulos y perfiles}

El análisis siguiente registra capacidades candidatas para la evolución del producto.
Para mantener la trazabilidad del Sprint 1, el alcance comprometido de RF-01 sigue
siendo la subida de bocetos, junto con su flujo y diseño de interfaz. Los demás
requerimientos y módulos requieren priorización y estimación propias antes de
considerarse parte de un sprint.

\subsubsection{Perfiles de usuario (roles)}

\begin{itemize}
    \item \textbf{Evaluador o investigador UX:} configura planes de prueba, modera sesiones, registra tiempos y documenta observaciones empíricas.
    \item \textbf{Líder o equipo de proyecto:} gestiona el Product Backlog y prioriza los hallazgos. En el equipo actual, estas responsabilidades corresponden al Product Owner y al equipo Scrum.
    \item \textbf{Docente o revisor:} consulta y audita la matriz de usabilidad, la trazabilidad del código y las comparativas de diseño. Sus permisos específicos quedan por definir.
\end{itemize}

Los perfiles anteriores describen actores del producto. Las personas participantes
de una prueba de usabilidad deben modelarse por separado de los roles de acceso a la
plataforma.

\subsubsection{Módulos principales propuestos}

\begin{itemize}
    \item \textbf{Módulo de evaluaciones:} asistente de configuración del test y bitácora para registrar resultados (éxito o fracaso), tiempos y errores por tarea.
    \item \textbf{Dashboard analítico:} panel con métricas consolidadas, tendencias de severidad y tasas de completitud. Su viabilidad depende de registrar ejecuciones de tareas y definir sus métricas.
    \item \textbf{Módulo de inteligencia artificial (\texttt{AiModule}):} componente asíncrono para resumir observaciones y generar borradores de historias en JSON estructurado.
    \item \textbf{Módulo de gestión ágil (SCRUM):} capacidad candidata para administrar backlog, sprints y retrospectivas. Durante el Sprint 1, Trello es el tablero de seguimiento; no se compromete construir un tablero duplicado ni su integración.
    \item \textbf{Módulo de exportación y evidencia:} generación de reportes Markdown y PDF, y consulta comparativa de diseños. Queda como capacidad futura por priorizar.
\end{itemize}

\subsubsection{Requerimientos funcionales candidatos}

Los siguientes requerimientos proceden del análisis del Product Owner y deben
refinarse, asignarse a identificadores y priorizarse en el Product Backlog:

\begin{itemize}
    \item Permitir la creación de un plan de prueba con objetivos, perfiles de participantes y guion de tareas.
    \item Registrar para cada ejecución de tarea el resultado (éxito o fracaso), la duración y la cantidad de errores.
    \item Procesar observaciones de texto libre mediante el servicio de IA y proponer una clasificación de problemas por categorías UX.
    \item Convertir una recomendación validada en una historia de usuario dentro del backlog definido por el equipo.
    \item Exportar los hallazgos validados en formatos integrables, inicialmente Markdown y PDF.
\end{itemize}

\subsubsection{Revisión técnica y de alcance}

\begin{itemize}
    \item \textbf{Arquitectura:} las capacidades pueden organizarse en módulos del backend NestJS y consumirse desde el frontend React. Para el producto inicial conviene mantener límites modulares sin asumir una arquitectura de microservicios.
    \item \textbf{Persistencia:} SQL Server es adecuado para usuarios, pruebas, tareas, observaciones e historias; la capa de acceso a datos debe seleccionarse para el backend NestJS. Para soportar las métricas propuestas habrá que modelar las ejecuciones de tareas; para imágenes conviene almacenar el archivo fuera de la base relacional y guardar en SQL Server sus metadatos y ubicación.
    \item \textbf{Procesamiento de IA:} una integración asíncrona puede implementarse mediante servicios del backend NestJS y un trabajador en segundo plano. El proveedor Gemini, un servicio denominado \texttt{GeminiService}, los DTO con \texttt{class-validator} y las pruebas Jest son propuestas compatibles del aporte de García, pero deben confirmarse antes de considerarse decisiones oficiales. Toda respuesta JSON debe validarse y mapearse a estructuras tipadas antes de persistirla.
    \item \textbf{Gestión ágil:} mantener Trello como fuente de verdad del Sprint Backlog evita duplicar el estado. Una futura integración con otro tablero requiere una decisión explícita del equipo.
    \item \textbf{Exportación:} Markdown, PDF y la integración con incidencias de GitHub son capacidades futuras propuestas; no forman parte de las dos historias comprometidas para el Sprint 1.
\end{itemize}

\subsection{RNF-01: Soluciones UX para gestionar la latencia de IA}

Las respuestas de modelos de lenguaje pueden variar en duración. Las siguientes
soluciones son viables en una interfaz web y deben acompañar, no sustituir, la
medición del objetivo de respuesta de menos de 10 segundos:

\begin{itemize}
    \item \textbf{Skeleton screens (carga progresiva):} mostrar una estructura aproximada de la historia mientras se procesa la solicitud. El esqueleto debe anunciarse accesiblemente y desaparecer al mostrar el resultado o un estado de error.
    \item \textbf{Microcopy de estado:} utilizar mensajes como ``Analizando observaciones'' o ``Preparando historias de usuario''. Mostrar fases concretas, como estructuración JSON, solo si el backend puede reportarlas; en caso contrario, usar un mensaje general sin simular progreso.
    \item \textbf{Efecto visual shimmer:} puede aplicarse de forma sutil a los contenedores de carga. Debe respetar contraste y movimiento reducido (\texttt{prefers-reduced-motion}); la animación es un complemento y no debe comunicar por sí sola el estado.
    \item \textbf{Degradación elegante:} el objetivo RNF-01 se incumple si el resultado no está disponible antes de los 10 segundos. La interfaz debe indicar que el análisis sigue pendiente y conservar el estado de la solicitud. El umbral de 15 segundos propuesto por el Product Owner puede utilizarse para activar una alternativa de espera o reintento, pero debe validarse y no debe presentarse como cumplimiento del RNF.
    \item \textbf{Fallback confiable:} ante una caída o error, informar al usuario, conservar la solicitud cuando sea posible y permitir reintentar sin crear duplicados. No se recomienda reemplazar el resultado real por una síntesis simulada; si se usa una respuesta simulada en una demostración, debe identificarse claramente como tal.
\end{itemize}

La medición del umbral debe acordar el perfil de prueba (ambiente, tamaño de imagen,
volumen de observaciones y punto inicial y final del cronómetro). La propuesta actual
considera desde la recepción de la solicitud hasta que las historias generadas están
disponibles para el evaluador; esta definición debe validarse antes de las pruebas.

\section{Historias de usuario y criterios de aceptación}

Los siguientes textos formalizan el alcance compartido para el Sprint 1. Las
estimaciones son propuestas para planificación, no compromisos definitivos: deben
validarse con el equipo. Los formatos, tamaños y límites de archivos, así como las
condiciones de medición del tiempo de respuesta, deben ser definidos por el Product
Owner y el Development Team.

\subsection{RF-01 -- Subida de bocetos}

\vspace{0.5cm}
\begin{longtable}{|p{0.45\textwidth}|p{0.45\textwidth}|}
\hline
\rowcolor{gray!20}
\multicolumn{2}{|l|}{\textbf{Historia de Usuario: RF-01 -- Subida de bocetos}} \\ \hline
\textbf{Número:} RF-01 & \textbf{Usuario:} Evaluador UX \\ \hline
\multicolumn{2}{|l|}{\textbf{Nombre historia:} Subida de bocetos} \\ \hline
\textbf{Prioridad de negocio:} Alta & \textbf{Riesgo en desarrollo:} Medio \\ \hline
\textbf{Puntos estimados:} 5 SP & \textbf{Iteración asignada:} Sprint 1 \\ \hline
\multicolumn{2}{|l|}{\textbf{Programador responsable:} Manolo y Ariel (Development Team)} \\ \hline
\multicolumn{2}{|p{0.9\textwidth}|}{
\textbf{Descripción:}\newline
COMO evaluador UX,\newline
QUIERO subir imágenes o bocetos de interfaces a la plataforma,\newline
PARA incorporarlos al análisis de usabilidad asistido por inteligencia artificial y generar historias de usuario para el Product Backlog.
} \\ \hline
\multicolumn{2}{|l|}{\textbf{Validación:}} \\ \hline
\multicolumn{2}{|p{0.9\textwidth}|}{
\textbf{Escenario 1: Selección y vista previa}\newline
DADO: Que el evaluador UX se encuentra en la pantalla de carga de una prueba de usabilidad,\newline
CUANDO: Selecciona una imagen o un boceto desde su dispositivo,\newline
ENTONCES: La plataforma muestra el archivo seleccionado y una vista previa, y permite confirmar o cancelar la carga.
} \\ \hline
\multicolumn{2}{|p{0.9\textwidth}|}{
\textbf{Escenario 2: Carga correcta}\newline
DADO: Que el archivo cumple las restricciones definidas por la plataforma,\newline
CUANDO: El evaluador confirma la carga,\newline
ENTONCES: La plataforma confirma el resultado, asocia el archivo a la prueba y lo deja disponible para el análisis.
} \\ \hline
\end{longtable}
\vspace{0.5cm}

\subsubsection*{Tareas y evidencia de diseño}

\begin{table}[htbp]
    \centering
    \caption{Desglose inicial de trabajo para RF-01}
    \label{tab:tareas-rf01}
    \begin{tabularx}{\textwidth}{L{1.7cm}Y L{2.7cm} L{1.7cm}}
        \toprule
        \textbf{ID} & \textbf{Tarea} & \textbf{Responsable} & \textbf{Estado} \\
        \midrule
        RF-01.1 & Analizar los requerimientos base de carga & Luis; revisión del Development Team & En progreso \\
        RF-01.2 & Definir el flujo de usuario para cargar un boceto & Luis; revisión del Development Team & En progreso \\
        RF-01.3 & Diseñar las pantallas de carga & Luis; revisión del Development Team & En progreso \\
        \bottomrule
    \end{tabularx}
\end{table}

\subsection{RNF-01 -- Tiempo de respuesta de la IA menor a 10 segundos}

\vspace{0.5cm}
\begin{longtable}{|p{0.45\textwidth}|p{0.45\textwidth}|}
\hline
\rowcolor{gray!20}
\multicolumn{2}{|l|}{\textbf{Historia de Usuario: RNF-01 -- Tiempo de respuesta de IA}} \\ \hline
\textbf{Número:} RNF-01 & \textbf{Usuario:} Evaluador UX \\ \hline
\multicolumn{2}{|l|}{\textbf{Nombre historia:} Tiempo de respuesta de IA menor a 10 segundos} \\ \hline
\textbf{Prioridad de negocio:} Por definir con el PO & \textbf{Riesgo en desarrollo:} Por evaluar \\ \hline
\textbf{Puntos estimados:} 8 SP & \textbf{Iteración asignada:} Sprint 1 \\ \hline
\multicolumn{2}{|l|}{\textbf{Programador responsable:} Manolo y Ariel (Development Team)} \\ \hline
\multicolumn{2}{|p{0.9\textwidth}|}{
\textbf{Descripción:}\newline
COMO evaluador UX,\newline
QUIERO recibir el resultado del análisis de inteligencia artificial en menos de 10 segundos y conocer su estado mientras se procesa,\newline
PARA revisar las observaciones y las historias de usuario generadas sin esperas prolongadas.
} \\ \hline
\multicolumn{2}{|l|}{\textbf{Validación:}} \\ \hline
\multicolumn{2}{|p{0.9\textwidth}|}{
\textbf{Escenario 1: Análisis dentro del umbral objetivo}\newline
DADO: Que el evaluador inicia un análisis con una imagen y las observaciones requeridas, bajo un perfil de prueba previamente acordado,\newline
CUANDO: El sistema recibe la solicitud de análisis,\newline
ENTONCES: Las historias de usuario generadas quedan disponibles en menos de 10 segundos desde la recepción de la solicitud.
} \\ \hline
\multicolumn{2}{|p{0.9\textwidth}|}{
\textbf{Escenario 2: Retroalimentación durante el procesamiento}\newline
DADO: Que el evaluador ha iniciado un análisis y el resultado todavía no está disponible,\newline
CUANDO: El sistema procesa la solicitud,\newline
ENTONCES: La interfaz comunica que el análisis está en curso mediante un indicador de carga o un esqueleto, según la solución UX definida.
} \\ \hline
\multicolumn{2}{|p{0.9\textwidth}|}{
\textbf{Escenario 3: El análisis supera el umbral objetivo}\newline
DADO: Que el resultado no está disponible al cumplirse 10 segundos,\newline
CUANDO: El procesamiento continúa,\newline
ENTONCES: La interfaz comunica que existe una demora, no muestra la operación como completada y registra el incumplimiento para validación técnica.
} \\ \hline
\end{longtable}
\vspace{0.5cm}

\subsubsection*{Trabajo previsto}

\begin{itemize}
    \item Investigar patrones UX para comunicar latencia, incluidos indicadores de carga y esqueletos.
    \item Definir el enfoque para medir y validar el umbral de respuesta de menos de 10 segundos.
    \item Documentar la arquitectura oficial del proyecto, estableciendo explícitamente que el Backend funciona con NestJS y el Frontend con React.
\end{itemize}

\section{Avances y productos del Sprint 1}

Al momento de elaborar este documento, se reportan los siguientes avances:

\begin{itemize}
    \item Repositorio del proyecto creado en GitHub.
    \item Tablero Kanban de Trello estructurado para gestionar el trabajo.
    \item Script inicial del modelo entidad-relación para SQL Server finalizado.
    \item RF-01 y RNF-01 en progreso.
\end{itemize}

Los resultados funcionales y el incremento verificable del producto se completarán
con evidencia al cierre del Sprint 1.

\section{Evidencias de UX/UI}

La matriz de evaluación solicita evidencias de diseño y análisis UX/UI. Se registrará
en esta sección el artefacto correspondiente y su ubicación o enlace cuando esté
disponible.

\subsection{User Persona}
Evidencia del perfil de usuario principal de la plataforma, detallando su contexto,
objetivos y frustraciones.
\begin{center}
    \includegraphics[width=\textwidth]{imagenes/User_Persona.png}
\end{center}

\subsection{Journey Map}
Mapeo del viaje del evaluador durante la configuración de la prueba y la carga de
bocetos, identificando los puntos de fricción y las oportunidades UX para gestionar
la latencia.
\begin{center}
    \includegraphics[width=\textwidth]{imagenes/Journey_Map.png}
\end{center}

\subsection{User Flow de carga de bocetos}
\begin{center}
    \includegraphics[width=0.9\textwidth]{imagenes/userFlow.png}
\end{center}

\subsection{Mockups de carga}
Estos wireframes de baja fidelidad documentan los estados interactivos de reposo,
loading y resolución, diseñados para mitigar la latencia de la inteligencia artificial
(RNF-01).
\begin{center}
    \includegraphics[width=0.9\textwidth]{imagenes/imgGarcia1.png}
\end{center}
\begin{center}
    \includegraphics[width=0.9\textwidth]{imagenes/imgGarcia2.png}
\end{center}
\begin{center}
    \includegraphics[width=0.9\textwidth]{imagenes/imgGarcia3.png}
\end{center}

\subsubsection{Wireframes y mockups de interfaz de Ariel}

Las siguientes evidencias presentan los wireframes y mockups elaborados por Ariel para
RF-01. En conjunto, definen la jerarquía visual de la pantalla y la estructura de sus
componentes, desde la selección del archivo hasta la confirmación de la carga.

Estos diseños sirven como referencia para la implementación del frontend en React y
representan los elementos y estados del flujo de subida de bocetos. Su organización
visual permite al evaluador comprender las acciones disponibles y el resultado de
cada paso de la interacción.

\begin{center}
    \includegraphics[width=0.9\textwidth]{imagenes/imgArielBajaFidelidad.png}
\end{center}
\vspace{0.5cm}
\begin{center}
    \includegraphics[width=0.9\textwidth]{imagenes/imgAriel2.png}
\end{center}
\vspace{0.5cm}
\begin{center}
    \includegraphics[width=0.9\textwidth]{imagenes/imgAriel3.png}
\end{center}
\vspace{0.5cm}
\begin{center}
    \includegraphics[width=0.9\textwidth]{imagenes/imgAriel4.png}
\end{center}
\vspace{0.5cm}
\begin{center}
    \includegraphics[width=0.9\textwidth]{imagenes/imgAriel5.png}
\end{center}

\section{Base técnica y modelo de datos}

\subsection{Repositorio y código base}
El repositorio de GitHub contiene la estructura inicial del proyecto y el script DDL
de la base de datos relacional para SQL Server, junto con este documento LaTeX. La
siguiente evidencia muestra los commits correspondientes.
\begin{center}
    \includegraphics[width=\textwidth]{imagenes/repositorio_github.png}
\end{center}

\subsection{Base de datos}
Se finalizó un script inicial del modelo entidad-relación para SQL Server. Se debe
incorporar el diagrama y validar el modelo frente a las entidades del proyecto:
usuarios, pruebas, tareas, observaciones, historias de usuario y bocetos. La tabla
\texttt{Bocetos} almacena los metadatos y la referencia de almacenamiento de cada
imagen, vinculada mediante llave foránea a \texttt{PruebasUsabilidad}. El archivo
binario se conserva fuera de SQL Server. El registro detallado de las ejecuciones de
tareas para las métricas analíticas queda pendiente de modelado.

\subsection{Stack backend y frontend}
El documento de arquitectura del proyecto establece explícitamente que el Backend
funciona con NestJS y el Frontend con React. Esta definición oficial orienta el
desarrollo de las próximas iteraciones y se alinea con las herramientas del equipo de
desarrollo, descartando cualquier alternativa técnica tentativa evaluada previamente.
La persistencia relacional utiliza SQL Server. El uso de Google Gemini, un servicio
\texttt{GeminiService}, validación mediante \texttt{class-validator} y pruebas con
Jest aparece en el aporte de García como una propuesta compatible, pero requiere
aprobación antes de registrarse como decisión oficial.

\section{Testing y validación}

La validación de los criterios de aceptación de RF-01 y RNF-01 está pendiente. Al
ejecutar las pruebas, se documentarán el caso, los datos utilizados, el resultado
esperado, el resultado obtenido y la evidencia correspondiente. Para RNF-01 también
se registrará el tiempo observado y el método de medición acordado. El perfil de
prueba de rendimiento (ambiente, tamaño de imagen y volumen de observaciones) debe
quedar documentado antes de evaluar el umbral de 10 segundos.

\section{Gestión ágil del Sprint}

\subsection{Tablero de Trello}
El Sprint Backlog se gestionó mediante Trello como herramienta visual para organizar
las historias de usuario y sus tareas, asignar responsables y dar seguimiento a los
estados del trabajo.
\begin{center}
    \includegraphics[width=\textwidth]{imagenes/tablero_kanban.png}
\end{center}

El aporte de García registra tarjetas de tareas para RF-01 y RNF-01; estas son tareas
de las dos historias del Sprint 1, no historias adicionales. Los identificadores,
responsables abreviados (MG y CI) y estados indicados en ese aporte deben cotejarse con
Trello y homologarse con el desglose de tareas de este documento. La tarea de pipeline
de exportación PNG/PDF se excluye del alcance de RF-01 en este sprint y queda como
candidata para priorización futura.

\subsection{Ceremonias y seguimiento}
Kevin, como Scrum Master, facilita las Daily Standups y la planificación del sprint,
y da seguimiento a los bloqueos. Las fechas, acuerdos y acciones se registrarán como
evidencia de gestión ágil.

\section{Retrospectiva del Sprint}

\subsection{Lo que salió bien}
\begin{itemize}
    \item Se consolidaron a tiempo las evidencias UX/UI del Sprint 1: User Persona, Journey Map y wireframes.
    \item Se estructuró la base del repositorio y el modelo relacional inicial en SQL Server.
\end{itemize}

\subsection{Oportunidades de mejora}
\begin{itemize}
    \item Se produjo una desalineación inicial sobre el stack: se consideró .NET, aunque la arquitectura oficial establece NestJS para el backend y React para el frontend.
    \item El modelo DDL requiere una entidad adicional para almacenar los metadatos y la referencia de los archivos de bocetos.
\end{itemize}

\subsection{Acciones para el siguiente sprint}
\begin{itemize}
    \item Comunicar y registrar las decisiones oficiales de arquitectura desde el inicio del sprint.
    \item Priorizar la actualización del script SQL Server para incorporar la entidad de archivos o bocetos y su relación con la prueba correspondiente.
\end{itemize}

\section{Métricas y conclusiones}

\subsection{Métricas del Sprint 1}
\begin{itemize}
    \item \textbf{Puntos de historia completados:} 13 SP de 13 SP planificados.
    \item \textbf{RF-01 -- Subida de bocetos:} 5 SP completados.
    \item \textbf{RNF-01 -- Latencia de IA menor a 10 segundos:} 8 SP completados.
    \item \textbf{Capacidad del equipo:} 32 horas planificadas y consumidas.
\end{itemize}

\subsection{Conclusiones}
\begin{itemize}
    \item Se cumplió al 100\% la Definition of Done del Sprint 1 definida por la matriz de evaluación.
    \item Se entregaron las bases de UX/UI, requerimientos y arquitectura necesarias para iniciar la fase de codificación en el Sprint 2.
\end{itemize}

\end{document}
